\section{Conclusions}
\label{sec:conclusions}

We revisited a real-world process selection and lot-sizing problem
whose original plant order book was solved by MIP, but whose enlarged
benchmark instances motivated heuristic development. The study
confronted a modern MIP solver, temporal-decomposition matheuristics,
and a strong hybrid metaheuristic baseline using comparison criteria
fixed in advance, on a benchmark extended across seven derived instance
families.

Three answers emerge. First, on the practical status of the
problem: at the scale of the plant's order book it is simply solved,
with proven optima in seconds, and this remains true for the smallest
derived families. The first-line recommendation for practitioners is
therefore to try the monolith with a modern solver before building
heuristics. What decides difficulty at a given size, however, is not
only the size. If decision makers value inventory reduction, the
weighted formulation offers substantial inventory reductions, but these
are accompanied by higher shortage in seven of the ten 2X instances,
the only set in which both variants are solved to proven optimality;
the same direction holds descriptively in 31 of 40 instances across
2X--5X. Whether this trade-off is operationally acceptable requires
decision-maker assessment. Second, on the question of paradigms: when
the monolith does falter, what replaces it is not the pure
metaheuristic. On the scales where it was run, our strong hybrid
GRASP-ILS baseline is credited with the best of ten attempts within the
stated cutoff, but is still dominated by MIP-based methods. The
effective replacement is decomposition that keeps the solver in the
loop.
Third, on the question of \emph{when}: the method-choice map has
two genuine axes. In the tested 8X and 10X families, decomposition
wins the majority of instances at the 3,600-second cutoff; granting the
monolith triple time restores its lead in seven of those ten cases,
though decomposition remains superior in the other three. This
cross-budget comparison is asymmetric, because decomposition was not
run at 10,800 seconds, but it marks the practical frontier between
rolling re-planning deadlines and longer offline runs.

Beyond the verdicts, the study contributes methodological assets:
a fully specified, reproducible implementation of relax-and-fix and
fix-and-optimize in an algebraic modeling environment, including
the budget-guard and construction-floor policies whose absence we
showed can silently invalidate comparisons; a 60-instance derived
benchmark with best-known solutions, dual bounds, trajectories, and
per-window logs; and a worked example of stating decision criteria
before experimentation in computational OR, including two hypotheses
that were not supported and whose diagnosis improved the methods.

Limitations bound the claims. Our comparison with the original
study is a statement about practice, not a controlled measurement of
solver progress: solver version, hardware, parallelism, and objective
all differ, and disentangling them would require machine-independent
protocols we did not attempt. All benchmark instances derive from a
single industrial case, with only five base patterns represented in
the new 8X and 10X families, so the results are not a population sample
of grain-fusion order books. The MFP model prices setups implicitly
through the one-process-per-period structure rather than explicitly;
horizon replication, while faithful to the inherited benchmark, yields
periodic demand; matheuristic construction at
extreme scale fell back on a greedy floor, leaving the
relax-and-fix window policy suboptimal precisely where
decomposition matters most; decompositions were not run at the
largest budget; and the pure metaheuristic was not run on the two
largest families, so its comparison is bounded by $|T|\le 95$.
Future work follows directly: the shortage- and
setup-cost extensions of the original study under modern solvers;
stochastic order books; and, most immediately, replacing fixed
window policies with learned, state-dependent ones, the released
per-window logs, containing dual signals, incumbent profiles, and
realized improvements for every solve of this campaign, were
designed as training data for exactly that line of research.
